\section{Auth0 Acquisition: Primitive compartidas, sin prisa por fusionar todos los límites}

McKinnon describe a Okta y Auth0 como complementarias: la primera se ha orientado desde hace tiempo a workforce/IT, y la segunda es conocida por su customer identity developer-first. Ambas comparten las primitive de directory, authentication, authorization, protocol, security y scale, pero atienden a distintos buyer, end user, SDK, customization y release expectation. El primer principio de la arquitectura de una adquisición no es dibujar cuanto antes un solo org chart, sino proteger el customer contract.

\subsection{Workforce Identity y Customer Identity}

La workforce identity se ocupa de empleados, contractor, IT administration, enterprise app y lifecycle; la customer identity se ocupa de los usuarios finales del producto, developer integration, custom UX, social login, application-specific profile y tráfico de internet. Una plataforma común puede reutilizar protocol, threat intelligence y reliability practice, pero la runtime migration, el tenant model y la developer experience requieren validación independiente.

\lecturefigure{14-two-platform-acquisition.png}{Una adquisición debe respetar los límites de usuarios, flujos de trabajo e infraestructura que distinguen a la workforce identity de la customer identity.}{Subtítulos locales 30:00--36:20; materiales oficiales de la adquisición de Okta/Auth0 de 2021; redibujo conceptual.}

\begin{knowledgebox}{Lectura de la figura: una shared primitive no equivale a un shared runtime}
Ambos tipos de producto necesitan directory, OAuth/OIDC, MFA, security y scale; pero el workforce buyer puede exigir IT policy, employee lifecycle y enterprise integration, mientras que la customer identity pone el acento en SDK, experiencias con la marca propia, picos extremos y developer control. Después de la adquisición conviene unificar primero la threat intelligence, la support escalation y la strategic interface, y luego decidir con evidencia de compatibility, SLO y migration si se fusiona el runtime.
\end{knowledgebox}

\teachervoice{Al recordar la adquisición, el ponente considera que el escepticismo de los inversionistas hacia las M\&A fue mayor de lo que él esperaba; también reconoce que, en un entorno de tasas de interés bajas y crecimiento acelerado, parte de la integración del negocio pudo avanzar demasiado rápido. Esta reflexión advierte a los equipos de sistemas que el crecimiento de corto plazo oculta el integration cost, y que un viento macroeconómico a favor no debe confundirse con una arquitectura ya validada.}

\begin{warningbox}{Tres tipos de costos ocultos de la platform consolidation}
Unificar el empaquetado comercial puede ocurrir antes de que entitlement y billing sean coherentes; unificar la console puede ocultar distintos tenant/security model; unificar el backend amplía el blast radius, obliga a migrar estado y modifica los límites de cumplimiento normativo del cliente. Cada paso requiere rollback, data migration validation, customer communication y parallel run, en lugar de sustituir la aceptación técnica por «one company».
\end{warningbox}

\subsection{Resumen de la sección}

La integración posterior a una adquisición es una migración conjunta del product contract, del runtime state y de los incentivos de la organización. Compartir identity primitive puede generar sinergias, pero preservar los límites (preserve boundary) suele ser más seguro que unificar de inmediato.
